\section{Introduction}

The relationship between media and government has been extensively studied in the fields of economics and political science, and has been identified as a critical component of electoral accountability in Latin America \citep{Audits-Ferraz-QJE, PolInfo-Enriquez-JEEA, Marshall}. A well-functioning media wields significant power to hold government actions accountable and inform the public, thereby fostering transparency and responsiveness \citep{Stromberg-MediaCoverage, Account-Boix-JoLEO}. 

Crucially, incumbent governments are fully aware of the impact that media and the broader information environment have on their chances of remaining in power, prompting them to devote substantial resources to control the information reaching the public \citep{Autocracy-Guriev-JPE}. In Mexico, the decisions media outlets make—what stories to cover, how to cover them, and when—have been deeply tied to the government’s use of state resources to influence narratives, suppress dissent, and maintain political dominance. From outright authoritarian control during the 20th century \citep{FourthState-Lawson-Book} to subtler forms of influence through advertising spending \citep{Times, Contralinea-Desvio}, the relationship between the media and the state remains a key factor in shaping public discourse.

This study wishes to causally answer how the government reacts to the media being critical of it. The political science and economic literature have documented this censorship behavior in authoritarian contexts \citep{CensorshipCollective-King-APSR, CensorshipChina-King-Science}, and there is some historical and observational evidence on well functioning western democracies \citep{PressCriticism-Carlson-Journalism}, and on the Latin American context \citep{Peru-Macmillan-JEP, FourthState-Lawson-Book, ArgCorr-DiTella-AEJAE}. However, to the best of my knowledge, there is still no causal evidence of this behavior.

To answer if the government reacts to the media being critical, an interesting outcome at the media outlet level could be government spending in advertising. Since a simple OLS estimation using advertising spending by media outlets as the dependent variable and the number of critical articles they publish would likely suffer from endogeneity and unobserved confounding factors, I propose an instrumental variable approach. Specifically, I instrument the number of critical articles using exogenous variation in the assignment of seating at AMLO's morning press conferences. Starting in April 2022, the President and the Office of Communication overseeing the conferences began randomly assigning attending journalists to the first, second, or third row. 

To validate this instrument, I first demonstrate that being seated in the first row significantly increased the likelihood of a journalist asking a question in a sample of 10 randomly selected conferences. Additionally, I show that the distribution of journalists asking questions shifted after the lottery system was introduced—from one skewed in favor of pro-government media to a more balanced distribution, where critical outlets had greater opportunities to participate. This supports the argument that the seating assignment was effectively randomized.

My hypothesis then is that when media outlets write critical (favorable) coverage of the government, they get punished (compensated) by a reduction (increase) in the advertising spending in subsequent periods. Anecdotal evidence suggests that in past administrations the Mexican government has used this spending to silence critics and censor newspaper outlets \citep{Times}. Moreover, AMLO reduced this spending by approximately 78\% from the past administration \citep{Article19}, so now resources are more scarce. Given the substantial reduction, and the past government-media client-business like relationship, one could make the argument that through this channel the government can exert greater control of the media outlets as now more than ever they are more dependent on it to survive.

The rest of the research proposal is organized as follows. Section \ref{sec:lit_review} provides a detailed description of the literature related to this topic and the historical and political Mexican context. Section \ref{sec:data} outlines a description of all the data used in the analysis, as well as an explanation on how I will build each variable. Section \ref{sec:emp_strategy} presents the empirical strategy. Section \ref{sec:results} presents preliminary results.
